% !Mode:: "TeX:UTF-8"

\chapter{面向错误缓解的策略研究与实现}

\section{引言}

上一章围绕大语言模型推理过程中的错误识别问题，讨论了如何对模型生成的中间过程进行检查。
仅能识别错误并不足以提升模型的实际推理能力。对于复杂问答任务，模型往往需要先理解问题，
再决定是否调用自身知识进行检索式思考，随后整合生成的信息并给出最终答案。上述任一环节出现偏差，
都可能在后续推理中被继续放大。因此，面向错误缓解的关键问题在于：如何将错误识别结果转化为可学习的监督信号，
使模型在之后的生成过程中主动减少类似错误。

现有自检索强化学习方法通常以最终答案正确性作为奖励信号，对模型的完整输出进行结果级优化。
这类方法能够鼓励模型生成更高准确率的答案，但奖励信号较为稀疏，无法直接指出推理链中哪一步出现了错误，
也难以说明应如何修正错误。过程奖励模型（Process Reward Model, PRM）试图对中间步骤进行评分，
但在知识调用场景中，错误并不总是表现为单个步骤的局部错误。模型可能生成了不合适的查询，
也可能在看似合理的信息片段中混入事实偏差，还可能在整合信息时得到与证据不一致的结论。
这些现象使得单纯的过程打分难以提供稳定、细粒度且可操作的监督信号。

针对上述问题，本文提出一种基于反思验证数据的错误缓解策略。该策略首先收集自检索强化学习
（Self-Search Reinforcement Learning, SSRL）过程中的 rollout 数据，使用强验证模型对模型的
“思考-检索-回答”过程进行错误标注，并生成包含错误解释和修正建议的反思反馈。
随后，本文使用构造得到的反思验证数据对 Qwen2.5-3B-Instruct 模型进行监督微调，
再以微调后的模型作为初始策略继续进行强化学习。该方法的核心思想是先让模型学习“如何发现并解释自身知识调用错误”，
再利用结果奖励进一步优化最终回答表现，从而弥补仅依赖结果奖励时监督信号过于稀疏的问题。

\section{背景知识和相关工作}

\subsection{基于结果奖励的推理强化学习}

强化学习已被广泛用于提升大语言模型在复杂推理任务中的表现。与监督微调直接模仿给定答案不同，
强化学习根据模型输出的奖励对策略进行更新，更适合处理答案空间较大、推理路径多样的问题。
在问答任务中，常见奖励来自最终答案是否正确，此类奖励也被称为结果奖励（Outcome Reward）。
结果奖励具有实现简单、评价目标明确的优点，但只能在完整回答生成后给出反馈，
无法直接约束模型的中间推理过程。

组相对策略优化（Group Relative Policy Optimization, GRPO）是一类适用于大语言模型推理优化的强化学习方法。
其基本思想是对同一问题采样多条候选输出，并根据同组样本之间的相对奖励估计优势，
从而降低额外价值模型的依赖。在自检索问答任务中，模型不仅需要给出答案，
还需要在必要时生成检索式查询与相应信息，因此强化学习优化的对象是完整的推理与回答轨迹。
SSRL 方法正是利用这类结果奖励信号，鼓励模型更有效地调用自身参数化知识。

GRPO 通过对同一提示词生成多个输出，并在这些输出构成的群体内部计算相对奖励分数，以此来计算优势函数，生成奖励信号。具体而言，对于给定的问题 $q$，旧策略模型 $\pi_{\theta_{\mathrm{old}}}$ 会采样生成一个包含 $n$ 个候选回答的群体 $\{o_1,o_2,\cdots,o_n\}$。随后，奖励模型（或规则化奖励函数）会对每个回答 $o_i$ 进行打分得到 $r_i$。GRPO 通过在群体内部对奖励进行标准化处理来计算相对优势 $A_i$：
\begin{align}
A_i=\frac{r_i-\operatorname{mean}(r_1,\cdots,r_n)}
{\operatorname{std}(r_1,\cdots,r_n)}
\end{align}

基于当前策略 $\pi_\theta$ 与参考模型 $\pi_{\mathrm{ref}}$ 之间的 KL 散度惩罚项与相对优势 $A_i$，可获得用于梯度更新的优势函数 $\mathcal{J}_{GRPO}(\theta)$：

\begin{align}
r_{i,t}(\theta)
&=
\frac{\pi_\theta(o_{i,t}\mid q,o_{i,<t})}
{\pi_{\theta_{\mathrm{old}}}(o_{i,t}\mid q,o_{i,<t})}
\end{align}

\begin{align}
\ell_{i,t}(\theta)
&=
\min\left[
r_{i,t}(\theta)A_i,\,
\operatorname{clip}\!\left(
r_{i,t}(\theta),
1-\varepsilon,
1+\varepsilon
\right)A_i
\right]
-\beta D_{\mathrm{KL}}\!\left[
\pi_\theta \middle\| \pi_{\mathrm{ref}}
\right]
\end{align}

\begin{align}
\mathcal{J}_{GRPO}(\theta)
&=
\frac{1}{n}\sum_{i=1}^{n}
\frac{1}{|o_i|}\sum_{t=1}^{|o_i|}
\ell_{i,t}(\theta)
\end{align}

然而，结果级强化学习仍存在明显局限。若某条轨迹的最终答案错误，模型无法直接得知错误来源：
错误可能来自问题理解，也可能来自知识调用、信息生成或最终归纳。若某条轨迹的答案正确，
其中也可能包含冗余、偶然正确或局部不可靠的中间过程。只依赖最终奖励容易导致模型学习到与答案相关但与过程可靠性无关的模式，
难以稳定缓解推理过程中的知识调用错误。

\subsection{过程监督与反思反馈}

为了弥补结果奖励的稀疏性，过程监督方法尝试对推理链中的中间步骤进行评价。
其代表做法是训练 PRM，对每一步推理是否合理给出分数，并在搜索或强化学习过程中使用该分数辅助决策。
过程监督能够更直接地约束推理路径，但构造成本较高\cite{ref15_he2025large}。

反思反馈提供了另一种形式的推理过程改进思路。与仅给出二值标签或步骤分数不同，反思反馈要求验证模型说明错误发生的位置、
错误原因以及可能的修正方向。对于知识调用任务而言，
反思反馈可以显式描述查询不充分、生成信息不准确、证据与答案不一致等问题，
为后续监督微调提供更具可学习性的目标。

\subsection{自检索问答任务}

自检索问答任务要求模型在回答问题时显式生成“思考-检索-回答”过程,这种方法被认为可以提升模型使用自身参数化知识的能力\cite{fan2025ssrl}。
当模型判断自身缺少直接回答所需知识时，会生成 \verb|<search>| 查询，
随后生成 \verb|<information>| 信息片段，并基于这些信息继续推理和回答。
需要强调的是，该过程并不从外部数据库返回真实文档，而是让模型通过自生成的检索形式组织其参数化知识。
因此，错误缓解的目标不是提升外部检索器召回率，而是提升模型对自身知识的调用、检查和整合能力。

本文实验使用 PopQA、MuSiQue、HotpotQA、BamboogleQA 和 NQ 五个问答数据集。
这些数据集覆盖了短事实问答、多跳推理问答以及需要组合证据的问题类型。
在此类任务中，模型若只学习最终答案，很容易忽视中间知识调用过程的可靠性；
而通过反思验证数据进行监督微调，可以使模型在进入强化学习之前获得更明确的错误修正先验。

\section{方法设计}

\subsection{总体框架}

本文提出的错误缓解框架包括四个阶段。第一阶段，使用 SSRL 模型在问答任务上生成 rollout 数据，
保留问题、模型完整输出和标准答案。第二阶段，使用强验证模型对每条样本进行过程级检查，
判断模型输出中是否存在错误，并生成结构化反思反馈。第三阶段，将验证结果整理为反思验证数据集，
对 Qwen2.5-3B-Instruct 进行监督微调，使模型学习错误定位、错误解释和修正建议的生成方式。
第四阶段，以反思微调后的模型作为初始策略，在相同问答任务上继续进行强化学习，
得到本文方法的最终模型。

该方法在强化学习之前加入反思能力注入步骤。反思能力为模型参数化知识提供错误识别和修改的过程信息，对强化学习轨迹生成的优化可能起到优化作用。
二者结合后，模型既能保持面向最终准确率的优化目标，也能在生成过程中更主动地规避知识调用错误。

\subsection{反思验证数据构建}

反思验证数据来自 SSRL 训练过程中产生的 rollout 样本。每条样本包含三个核心部分：
问题、模型完整输出以及标准答案。模型完整输出保留了思考、检索式查询、生成信息和最终回答等内容，
因此验证模型可以从过程而非仅从答案层面对样本进行分析。

本文使用 DeepSeek V4 Flash 作为验证模型，对样本进行自动标注。标注提示词要求验证模型首先审查
“思考-检索-回答”全过程，判断其中是否存在错误；若存在错误，则输出错误描述和反思反馈；
若不存在错误，则说明各步骤为何合理。标注结果采用严格 JSON 格式。

本文共标注 10000 条反思验证样本，按照 8:2 的比例划分训练集和测试集。
训练集用于监督微调，测试集用于检查模型是否学习到稳定的过程验证与反思反馈能力。
与直接使用最终答案进行监督不同，该数据集的目标输出包含对错误原因的解释和修正建议，
因而能够为后续强化学习提供更适合过程改进的初始策略。

\subsection{反思微调与强化学习}

本文以 Qwen2.5-3B-Instruct 为基座模型，使用反思验证训练集进行指令微调。
训练目标是让模型在给定问题、模型输出和标准答案时，能够生成结构化的错误分析结果。
通过这一过程，模型学习到三类能力：其一，判断完整推理过程是否与标准答案一致；
其二，将错误定位到知识调用、信息生成或答案整合等具体环节；
其三，生成能够指导后续回答修正的反思反馈。

强化学习阶段，本文将反思微调后的模型接入 SSRL 强化学习流程，记为本文方法。
对照方法为未经过反思验证微调的 SSRL。为了保证比较公平，本文方法与 SSRL 使用相同的问答任务、
评价指标和结果奖励设置。强化学习过程中，模型仍以最终答案正确性作为主要优化目标；
不同之处在于，本文方法的初始策略已经通过反思验证数据学习到过程级错误意识，
在结果奖励信号较稀疏的情况下，能够形成可靠的知识调用。

\section{实验设置}

\subsection{数据集与评价指标}

实验结合使用 Natural Questions 和 HotpotQA 的作为自检索 SSRL 强化学习
训练数据集，以确保通用问答与多跳问答的混合。其中，Natural Question 数据集
收集了事实类知识问答，用于模型在一般问答任务场景中训练。HotpotQA 数据集包含的
问题需要多次调用知识才能得出答案，要求模型进行多跳推理。

实验使用 PopQA、MuSiQue、HotpotQA、BamboogleQA 和 NQ 五个测试集。
PopQA 和 NQ 更偏向开放域事实问答，MuSiQue 与 HotpotQA 侧重多跳推理，
BamboogleQA 则包含需要组合知识才能回答的问题。五个数据集共同用于检验方法在不同问答类型上的泛化能力。

评价指标采用 TensorBoard 中记录的 \texttt{reward mean@1}，在图表中记为 Accuracy。
该指标表示单次生成答案在测试集上的平均奖励，可直接反映模型最终回答的正确性。
本文分别报告训练完成检查点（final checkpoint）和最佳检查点（best checkpoint）的结果。
其中，final checkpoint 对应训练完成后的模型；best checkpoint 遵循相关工作设置，
选取测试过程中表现最优的检查点，用于观察方法在训练过程中的最佳可达性能。

\subsection{对比方法}

本文比较两种方法。第一种为 SSRL，即直接使用自检索强化学习进行结果级优化的基线方法。
第二种为本文方法，即先使用反思验证数据对 Qwen2.5-3B-Instruct 进行监督微调，
再继续执行相同的 SSRL 强化学习流程。该设置能够隔离反思验证训练的作用：
若本文方法优于 SSRL，则说明过程级反思监督能够为后续结果级强化学习提供有益初始化。

\section{实验结果与分析}

\subsection{训练完成检查点结果}

训练完成检查点上的逐数据集准确率如表~\ref{tab:error-mitigation-final-main} 所示，
包含训练步数和差值的详细结果见附录表~\ref{tab:appendix-error-mitigation-final}。

\begin{table}[htbp]
\caption{训练完成检查点上的测试集准确率对比}
\label{tab:error-mitigation-final-main}
\vspace{0.5em}\centering\wuhao
\begin{tabular}{lcc}
\toprule
数据集 & SSRL & 本文方法 \\
\midrule
PopQA & 0.3124 & 0.3322 \\
MuSiQue & 0.1576 & 0.1646 \\
HotpotQA & 0.2872 & 0.2908 \\
BamboogleQA & 0.2728 & 0.2656 \\
NQ & 0.3142 & 0.2926 \\
\bottomrule
\end{tabular}
\end{table}

从结果看，本文方法在 PopQA、MuSiQue 和 HotpotQA 上优于 SSRL，
其中多跳推理任务 PopQA 的提升最明显，准确率由 0.3124 提高到 0.3322，提升 0.0198。这一定程度上
说明反思验证训练对事实问答和部分多跳推理任务均有帮助。



图~\ref{fig:error-mitigation-final} 进一步展示了 final checkpoint 上各数据集的对比。
可以看到，本文方法的收益主要集中在 PopQA、MuSiQue 和 HotpotQA。
这三个数据集的共同特点是模型需要在生成过程中组织或组合知识，
反思反馈中关于查询、信息生成和答案整合的监督能够直接作用于这类错误。


\begin{figure}[htbp]
\centering
\includegraphics[width=0.92\textwidth]{error_mitigation_final_checkpoint}
\caption{训练完成检查点上的测试集准确率对比}
\label{fig:error-mitigation-final}
\end{figure}

\subsection{最佳检查点结果}

训练过程中最佳检查点的逐数据集准确率如表~\ref{tab:error-mitigation-best-main} 所示，
包含最佳检查点步数和差值的详细结果见附录表~\ref{tab:appendix-error-mitigation-best}。

\begin{table}[htbp]
\caption{最佳检查点上的测试集准确率对比}
\label{tab:error-mitigation-best-main}
\vspace{0.5em}\centering\wuhao
\begin{tabular}{lcc}
\toprule
数据集 & SSRL & 本文方法 \\
\midrule
PopQA & 0.3268 & 0.3466 \\
MuSiQue & 0.1646 & 0.1702 \\
HotpotQA & 0.3032 & 0.2962 \\
BamboogleQA & 0.3337 & 0.3440 \\
NQ & 0.3142 & 0.2980 \\
\bottomrule
\end{tabular}
\end{table}

与 final checkpoint 相比，best checkpoint 更能反映方法在合适训练阶段可达到的上限。
结果显示，本文方法在 PopQA、MuSiQue 和 BamboogleQA 上取得更高最佳准确率，说明反思验证训练整体上能够提升强化学习过程中的最佳可达性能。

图~\ref{fig:error-mitigation-best} 展示了最佳检查点结果。
图中可见在多数数据集上，本文方法带来了最佳性能的改善，本文方法在训练过程中形成了更高的性能上界。在 HotpotQA 和 NQ 上，本文方法未带来性能
提升，说明反思验证训练的收益与任务类型相关。

\begin{figure}[htbp]
\centering
\includegraphics[width=0.92\textwidth]{error_mitigation_best_checkpoint}
\caption{最佳检查点上的测试集准确率对比}
\label{fig:error-mitigation-best}
\end{figure}

\subsection{结果讨论}

综合 final checkpoint 与 best checkpoint 的结果可以看出，反思验证训练对自检索强化学习具有促进作用，
但这种促进并非在所有数据集上都稳定成立。本文方法在多数数据集的最佳检查点上优于 SSRL，
说明过程级反思监督能够提升模型在强化学习中的可达性能。其原因在于，
反思数据把原本隐藏在最终奖励背后的错误原因显式化，使模型在进入强化学习之前已经具备一定的自我检查能力。
当强化学习继续以结果奖励优化最终答案时，这种自我检查能力可以帮助模型生成更可靠的查询、
更一致的信息片段以及更符合证据的回答。

同时，实验也暴露出两个限制。第一，反思验证数据由强模型自动标注得到，其质量依赖验证模型本身。
若验证模型对某些问题类型的错误判断不稳定，监督微调可能引入偏差。第二，final checkpoint 与 best checkpoint
之间存在差异，说明训练过程仍需要更细致的早停和模型选择策略。特别是在 NQ 等短事实问答数据集上，
反思式推理可能带来额外生成负担，导致最终答案准确率下降。因此，后续工作可进一步研究按任务类型自适应地决定是否启用反思过程，
并结合更细粒度的验证信号改进强化学习中的奖励设计。

\section{本章小结}

本章围绕大语言模型推理错误的缓解问题，提出了一种基于反思验证数据的自检索强化学习改进方法。
该方法从 SSRL rollout 中构造反思验证样本，使用 DeepSeek V4 Flash 对模型的“思考-检索-回答”过程进行错误分析，
并将得到的 10000 条样本按 8:2 划分后用于 Qwen2.5-3B-Instruct 的监督微调。
在此基础上，本文继续进行 SSRL 强化学习，得到反思增强的问答模型。

实验结果表明，本文方法在 PopQA、MuSiQue 和 BamboogleQA 等数据集上提升了最佳检查点准确率，
五个数据集的平均最佳准确率由 0.2885 提高到 0.2910。
这说明反思验证监督能够在一定程度上弥补结果奖励信号稀疏的问题，
为强化学习提供更具过程意识的初始策略。

% Local Variables:
% TeX-master: "../thesis"
% TeX-engine: xetex
% End:
